Deep learning has transformed organ segmentation, yet a persistent gap remains between benchmark performance and surgical planning deployment. We identify which 3D CT organ segmentation algorithms are ready for clinical deployment in surgical planning and what practical barriers remain. Three databases (PubMed via Entrez E-utilities, arXiv via the official arXiv API, and Semantic Scholar via the Graph API) were searched for studies published between 2015 and January 2026. We included studies presenting deep learning methods for 3D organ segmentation from CT with quantitative evaluation on public benchmarks. A narrative synthesis of 52 included studies with benchmark-stratified performance comparison was conducted; formal meta-analysis was precluded by protocol heterogeneity. Auto-configured CNNs (nnU-Net, 82.3\% Dice) and pre-trained transformers (Swin UNETR, 82.4\%) achieve comparable leading accuracy, but domain shift, vascular segmentation deficits (65 to 75\% vs.\ 90\%+ for parenchymal organs), and integration complexity remain primary barriers to clinical adoption. CNNs dominate (98.1\%), with 3D U-Net variants appearing in 82.7\% of papers. Auto-configured CNNs achieve 85.2$\pm$2.1\% mean Dice on the BTCV multi-organ benchmark. Limitations include an English-only search, a 39.1\% full-text retrieval rate, and benchmark performance likely overestimating clinical accuracy by 5--15\%. We conclude that implementation quality matters more than architectural novelty. A phased adoption strategy starting with TotalSegmentator or nnU-Net is recommended for surgical planning platforms. The review protocol is available at \href{https://github.com/gameguild-gg/docdo-paper}{\faGithub\ gameguild-gg/docdo-paper}.
